% Options for packages loaded elsewhere
\PassOptionsToPackage{unicode}{hyperref}
\PassOptionsToPackage{hyphens}{url}
\documentclass[
]{article}
\usepackage{xcolor}
\usepackage[margin=1in]{geometry}
\usepackage{amsmath,amssymb}
\setcounter{secnumdepth}{-\maxdimen} % remove section numbering
\usepackage{iftex}
\ifPDFTeX
  \usepackage[T1]{fontenc}
  \usepackage[utf8]{inputenc}
  \usepackage{textcomp} % provide euro and other symbols
\else % if luatex or xetex
  \usepackage{unicode-math} % this also loads fontspec
  \defaultfontfeatures{Scale=MatchLowercase}
  \defaultfontfeatures[\rmfamily]{Ligatures=TeX,Scale=1}
\fi
\usepackage{lmodern}
\ifPDFTeX\else
  % xetex/luatex font selection
\fi
% Use upquote if available, for straight quotes in verbatim environments
\IfFileExists{upquote.sty}{\usepackage{upquote}}{}
\IfFileExists{microtype.sty}{% use microtype if available
  \usepackage[]{microtype}
  \UseMicrotypeSet[protrusion]{basicmath} % disable protrusion for tt fonts
}{}
\makeatletter
\@ifundefined{KOMAClassName}{% if non-KOMA class
  \IfFileExists{parskip.sty}{%
    \usepackage{parskip}
  }{% else
    \setlength{\parindent}{0pt}
    \setlength{\parskip}{6pt plus 2pt minus 1pt}}
}{% if KOMA class
  \KOMAoptions{parskip=half}}
\makeatother
\usepackage{color}
\usepackage{fancyvrb}
\newcommand{\VerbBar}{|}
\newcommand{\VERB}{\Verb[commandchars=\\\{\}]}
\DefineVerbatimEnvironment{Highlighting}{Verbatim}{commandchars=\\\{\}}
% Add ',fontsize=\small' for more characters per line
\usepackage{framed}
\definecolor{shadecolor}{RGB}{248,248,248}
\newenvironment{Shaded}{\begin{snugshade}}{\end{snugshade}}
\newcommand{\AlertTok}[1]{\textcolor[rgb]{0.94,0.16,0.16}{#1}}
\newcommand{\AnnotationTok}[1]{\textcolor[rgb]{0.56,0.35,0.01}{\textbf{\textit{#1}}}}
\newcommand{\AttributeTok}[1]{\textcolor[rgb]{0.13,0.29,0.53}{#1}}
\newcommand{\BaseNTok}[1]{\textcolor[rgb]{0.00,0.00,0.81}{#1}}
\newcommand{\BuiltInTok}[1]{#1}
\newcommand{\CharTok}[1]{\textcolor[rgb]{0.31,0.60,0.02}{#1}}
\newcommand{\CommentTok}[1]{\textcolor[rgb]{0.56,0.35,0.01}{\textit{#1}}}
\newcommand{\CommentVarTok}[1]{\textcolor[rgb]{0.56,0.35,0.01}{\textbf{\textit{#1}}}}
\newcommand{\ConstantTok}[1]{\textcolor[rgb]{0.56,0.35,0.01}{#1}}
\newcommand{\ControlFlowTok}[1]{\textcolor[rgb]{0.13,0.29,0.53}{\textbf{#1}}}
\newcommand{\DataTypeTok}[1]{\textcolor[rgb]{0.13,0.29,0.53}{#1}}
\newcommand{\DecValTok}[1]{\textcolor[rgb]{0.00,0.00,0.81}{#1}}
\newcommand{\DocumentationTok}[1]{\textcolor[rgb]{0.56,0.35,0.01}{\textbf{\textit{#1}}}}
\newcommand{\ErrorTok}[1]{\textcolor[rgb]{0.64,0.00,0.00}{\textbf{#1}}}
\newcommand{\ExtensionTok}[1]{#1}
\newcommand{\FloatTok}[1]{\textcolor[rgb]{0.00,0.00,0.81}{#1}}
\newcommand{\FunctionTok}[1]{\textcolor[rgb]{0.13,0.29,0.53}{\textbf{#1}}}
\newcommand{\ImportTok}[1]{#1}
\newcommand{\InformationTok}[1]{\textcolor[rgb]{0.56,0.35,0.01}{\textbf{\textit{#1}}}}
\newcommand{\KeywordTok}[1]{\textcolor[rgb]{0.13,0.29,0.53}{\textbf{#1}}}
\newcommand{\NormalTok}[1]{#1}
\newcommand{\OperatorTok}[1]{\textcolor[rgb]{0.81,0.36,0.00}{\textbf{#1}}}
\newcommand{\OtherTok}[1]{\textcolor[rgb]{0.56,0.35,0.01}{#1}}
\newcommand{\PreprocessorTok}[1]{\textcolor[rgb]{0.56,0.35,0.01}{\textit{#1}}}
\newcommand{\RegionMarkerTok}[1]{#1}
\newcommand{\SpecialCharTok}[1]{\textcolor[rgb]{0.81,0.36,0.00}{\textbf{#1}}}
\newcommand{\SpecialStringTok}[1]{\textcolor[rgb]{0.31,0.60,0.02}{#1}}
\newcommand{\StringTok}[1]{\textcolor[rgb]{0.31,0.60,0.02}{#1}}
\newcommand{\VariableTok}[1]{\textcolor[rgb]{0.00,0.00,0.00}{#1}}
\newcommand{\VerbatimStringTok}[1]{\textcolor[rgb]{0.31,0.60,0.02}{#1}}
\newcommand{\WarningTok}[1]{\textcolor[rgb]{0.56,0.35,0.01}{\textbf{\textit{#1}}}}
\usepackage{graphicx}
\makeatletter
\newsavebox\pandoc@box
\newcommand*\pandocbounded[1]{% scales image to fit in text height/width
  \sbox\pandoc@box{#1}%
  \Gscale@div\@tempa{\textheight}{\dimexpr\ht\pandoc@box+\dp\pandoc@box\relax}%
  \Gscale@div\@tempb{\linewidth}{\wd\pandoc@box}%
  \ifdim\@tempb\p@<\@tempa\p@\let\@tempa\@tempb\fi% select the smaller of both
  \ifdim\@tempa\p@<\p@\scalebox{\@tempa}{\usebox\pandoc@box}%
  \else\usebox{\pandoc@box}%
  \fi%
}
% Set default figure placement to htbp
\def\fps@figure{htbp}
\makeatother
\setlength{\emergencystretch}{3em} % prevent overfull lines
\providecommand{\tightlist}{%
  \setlength{\itemsep}{0pt}\setlength{\parskip}{0pt}}
\usepackage{bookmark}
\IfFileExists{xurl.sty}{\usepackage{xurl}}{} % add URL line breaks if available
\urlstyle{same}
\hypersetup{
  pdftitle={Conjugate Priors},
  hidelinks,
  pdfcreator={LaTeX via pandoc}}

\title{Conjugate Priors}
\author{}
\date{\vspace{-2.5em}2026-04-17}

\begin{document}
\maketitle

\subsection{Introduction}\label{introduction}

Before MCMC made arbitrary Bayesian models tractable, conjugate priors
provided the only route to a tractable posterior. A prior is
\textbf{conjugate} to a likelihood if the resulting posterior belongs to
the same distributional family as the prior -- the prior and posterior
share a form, with updated parameters that have simple closed-form
expressions. Conjugate pairs remain pedagogically central because they
make the Bayesian update arithmetic, and they remain practically useful
as building blocks within larger hierarchical models.

\subsection{Prerequisites}\label{prerequisites}

The reader should know Bayes' theorem and be comfortable with the beta,
normal, and gamma distributions. Familiarity with basic R plotting is
assumed.

\subsection{Theory}\label{theory}

\subsubsection{Beta-Binomial}\label{beta-binomial}

For a binomial observation \(Y \sim \text{Binomial}(n, \theta)\) with
success probability \(\theta\) and a beta prior
\(\theta \sim \text{Beta}(\alpha, \beta)\), the posterior is

\[\theta \mid Y = y \;\sim\; \text{Beta}(\alpha + y, \beta + n - y).\]

The update rule is nothing more than ``add successes to the first shape
parameter and failures to the second''. The posterior mean is
\((\alpha + y)/(\alpha + \beta + n)\), a weighted average of the prior
mean \(\alpha/(\alpha + \beta)\) and the observed proportion \(y/n\),
with weights determined by the prior's ``effective sample size''
\(\alpha + \beta\).

\subsubsection{Normal-Normal with known
variance}\label{normal-normal-with-known-variance}

For observations
\(Y_1, \ldots, Y_n \overset{iid}{\sim} \mathcal{N}(\mu, \sigma^2)\) with
\(\sigma^2\) known, and a normal prior
\(\mu \sim \mathcal{N}(\mu_0, \tau_0^2)\), the posterior is

\[\mu \mid y \;\sim\; \mathcal{N}(\mu_n, \tau_n^2),\]

where the posterior precision is
\(1/\tau_n^2 = 1/\tau_0^2 + n/\sigma^2\), and the posterior mean is a
precision-weighted average
\(\mu_n = \tau_n^2 (\mu_0/\tau_0^2 + n \bar{y}/\sigma^2)\).

\subsubsection{Gamma-Poisson}\label{gamma-poisson}

For independent counts \(Y_i \sim \text{Poisson}(\lambda)\) and a gamma
prior \(\lambda \sim \text{Gamma}(\alpha, \beta)\) (shape-rate
parameterisation), the posterior is

\[\lambda \mid y \;\sim\; \text{Gamma}\!\left(\alpha + \sum y_i,\; \beta + n\right).\]

Each of these updates can be memorised as ``add the data to the prior
parameters'', with the precise form determined by the sufficient
statistics of the likelihood.

\subsection{Assumptions}\label{assumptions}

Conjugate analyses assume the same independence and distributional
conditions as the corresponding likelihoods. In addition, they assume
that the prior is genuinely an expression of prior belief, not a
mathematical convenience chosen for conjugacy alone. An informative
conjugate prior that contradicts the data produces a posterior that is a
compromise between the two; if this is not what you intended, the prior
was a bad choice.

\subsection{R Implementation}\label{r-implementation}

A visualisation of the beta-binomial update makes the idea concrete:

\begin{Shaded}
\begin{Highlighting}[]
\FunctionTok{library}\NormalTok{(ggplot2)}

\NormalTok{alpha\_0 }\OtherTok{\textless{}{-}} \DecValTok{2}
\NormalTok{beta\_0  }\OtherTok{\textless{}{-}} \DecValTok{2}
\NormalTok{y       }\OtherTok{\textless{}{-}} \DecValTok{8}
\NormalTok{n       }\OtherTok{\textless{}{-}} \DecValTok{10}

\NormalTok{alpha\_n }\OtherTok{\textless{}{-}}\NormalTok{ alpha\_0 }\SpecialCharTok{+}\NormalTok{ y}
\NormalTok{beta\_n  }\OtherTok{\textless{}{-}}\NormalTok{ beta\_0 }\SpecialCharTok{+}\NormalTok{ n }\SpecialCharTok{{-}}\NormalTok{ y}

\NormalTok{theta }\OtherTok{\textless{}{-}} \FunctionTok{seq}\NormalTok{(}\DecValTok{0}\NormalTok{, }\DecValTok{1}\NormalTok{, }\AttributeTok{length.out =} \DecValTok{401}\NormalTok{)}
\NormalTok{df }\OtherTok{\textless{}{-}} \FunctionTok{data.frame}\NormalTok{(}
  \AttributeTok{theta =} \FunctionTok{rep}\NormalTok{(theta, }\DecValTok{3}\NormalTok{),}
  \AttributeTok{density =} \FunctionTok{c}\NormalTok{(}\FunctionTok{dbeta}\NormalTok{(theta, alpha\_0, beta\_0),}
              \FunctionTok{dbinom}\NormalTok{(y, n, theta) }\SpecialCharTok{*}\NormalTok{ n,}
              \FunctionTok{dbeta}\NormalTok{(theta, alpha\_n, beta\_n)),}
  \AttributeTok{kind =} \FunctionTok{factor}\NormalTok{(}\FunctionTok{rep}\NormalTok{(}\FunctionTok{c}\NormalTok{(}\StringTok{"Prior"}\NormalTok{, }\StringTok{"Likelihood (scaled)"}\NormalTok{, }\StringTok{"Posterior"}\NormalTok{),}
                    \AttributeTok{each =} \DecValTok{401}\NormalTok{),}
                \AttributeTok{levels =} \FunctionTok{c}\NormalTok{(}\StringTok{"Prior"}\NormalTok{, }\StringTok{"Likelihood (scaled)"}\NormalTok{, }\StringTok{"Posterior"}\NormalTok{))}
\NormalTok{)}

\FunctionTok{ggplot}\NormalTok{(df, }\FunctionTok{aes}\NormalTok{(}\AttributeTok{x =}\NormalTok{ theta, }\AttributeTok{y =}\NormalTok{ density, }\AttributeTok{colour =}\NormalTok{ kind, }\AttributeTok{linetype =}\NormalTok{ kind)) }\SpecialCharTok{+}
  \FunctionTok{geom\_line}\NormalTok{(}\AttributeTok{linewidth =} \DecValTok{1}\NormalTok{) }\SpecialCharTok{+}
  \FunctionTok{labs}\NormalTok{(}\AttributeTok{x =} \FunctionTok{expression}\NormalTok{(theta), }\AttributeTok{y =} \StringTok{"Density"}\NormalTok{,}
       \AttributeTok{title =} \StringTok{"Beta{-}Binomial update with y = 8 of n = 10"}\NormalTok{) }\SpecialCharTok{+}
  \FunctionTok{theme\_minimal}\NormalTok{(}\AttributeTok{base\_size =} \DecValTok{12}\NormalTok{)}
\end{Highlighting}
\end{Shaded}

A \texttt{Beta(2,\ 2)} prior (symmetric, weakly informative around 0.5)
combined with 8 successes in 10 trials gives a \texttt{Beta(10,\ 4)}
posterior concentrated around 0.71.

\subsection{Output \& Results}\label{output-results}

Posterior summaries from the beta distribution are:

\begin{itemize}
\tightlist
\item
  Mean: \((\alpha + y)/(\alpha + \beta + n) = 10/14 \approx 0.714\)
\item
  95\% equal-tailed credible interval:
  \texttt{qbeta(c(0.025,\ 0.975),\ 10,\ 4)} \(\approx (0.47, 0.91)\)
\item
  Probability that \(\theta > 0.5\): \texttt{1\ -\ pbeta(0.5,\ 10,\ 4)}
  \(\approx 0.95\).
\end{itemize}

\subsection{Interpretation}\label{interpretation}

The posterior in a conjugate model is not an approximation; it is the
exact posterior under the stated prior. A Bayesian reporting this
analysis would write: ``With a weakly informative \(\text{Beta}(2, 2)\)
prior and data of 8 successes in 10 trials, the posterior for \(\theta\)
is \(\text{Beta}(10, 4)\), with posterior mean 0.71 (95\% CrI 0.47 to
0.91). The posterior probability that \(\theta\) exceeds 0.5 is 0.95.''

\subsection{Practical Tips}\label{practical-tips}

\begin{itemize}
\tightlist
\item
  Report the prior explicitly so readers can judge its influence.
\item
  Use informative priors when justified by domain knowledge; use weakly
  informative priors as the default for exploratory analyses.
\item
  Conjugate updating extends to multi-parameter cases
  (normal-inverse-gamma for \((\mu, \sigma^2)\), Dirichlet-multinomial
  for categorical probabilities), but outside the exponential family it
  usually does not hold.
\item
  For small-sample binomial problems the beta-binomial is particularly
  useful: it avoids the ad-hoc continuity corrections of frequentist
  proportion tests and gives natural credible intervals.
\item
  When your likelihood is non-conjugate, do not force a conjugate prior
  on it; use MCMC or variational methods instead.
\end{itemize}

\subsection{For Reviewers}\label{for-reviewers}

\textbf{What to look for in a paper using this method.}

\begin{itemize}
\tightlist
\item
  \textbf{Common misapplications.}
\item
  \textbf{Diagnostics that should be reported but often aren't.}
\item
  \textbf{Red flags in tables and figures.}
\item
  \textbf{What to verify.}
\item
  \textbf{What an adequate Methods paragraph must contain.}
\end{itemize}

\subsection{See also --- labs in this
chapter}\label{see-also-labs-in-this-chapter}

\begin{itemize}
\tightlist
\item
  \href{labs/lab-bayesian-thinking-control-vs-decision-error.Rmd}{Bayesian
  thinking; α control vs decision error}
\item
  \href{labs/lab-brms-stan-loo-hierarchical-models.Rmd}{brms / Stan,
  LOO, hierarchical models}
\end{itemize}

\end{document}
